\scanpage{631}
\mspage{619}

it is the subgroup \inserted{$\{1,\dot{\sigma}_{\infty}\}$} of $\mathcal{G}_3$.
(for the image of $\tau'\inserted{{}=\sigma\tau}\in\mathcal{M}_{0,3}$ in $\mathcal{G}_3$
\deleted{\unsure{is}} (since that of $\tau$ is $1$) is equal to that of
$\sigma=\sigma_{\infty}$, i.e. $\dot{\sigma}_{\infty}$. This means
that \deleted{\unsure{T}} we have
\[
T\simeq T_0\mathbin{\wedge}_{\{1,\dot{\sigma}_{\infty}\}}\mathcal{G}_3
\]
i.e. $T$ is obtained from a $\mu$-torsor (i.e. from a
quadratic extension of $\mathbb{Q}$\unsure{)} --- namely precisely
$\mathbb{Q}(\sqrt[3]{1})=\mathbb{Q}(j))$ by extension of the groups
of operators \inserted{$\mu\longrightarrow\mathcal{G}_3\quad(-1\longmapsto\dot{\sigma}_{\infty})$}.
Thus, taking the inverse image of $\mu=\{1,\dot{\sigma}_{\infty}\}$ in
$\widetilde{\mathcal{G}}_3$, we find an extension of $\mu$ by $\mu$
--- it can hardly be anything but $\mathbb{Z}/4\mathbb{Z}$ ! --- which I feel like denoting $\mathcal{E}''$

\[
\diagram{\begin{tikzcd}[ampersand replacement=\&,column sep=1.3em,row sep=0.15em]
\&\& \mathbb{Z}/4\mathbb{Z}\arrow[d,equal] \&\&\\
1\arrow[r]\&\mu\arrow[r]\&\mathcal{E}''\arrow[r]\&\mu\arrow[r,no head]\&1
\end{tikzcd}}
\]

\margin{yes, this is obvious since $\sigma=\sigma_{\infty}$ in
$\widetilde{\mathcal{G}}_3=\mathrm{SL}(2,\mathbb{Z})/\unsure{\Pi_0|\Pi}$
is indeed of order $4$, so $\mathcal{E}''\simeq\mathbb{Z}/4\mathbb{Z}$
\unsure{canonically}\ldots}
and the \unsure{problem} of obstruction is to find a lifting

\[
\diagram{\begin{tikzcd}[ampersand replacement=\&,column sep=3em,row sep=2em]
\mathcal{E}''\arrow[r,no head]\&\mu\\
\Gamma_{\mathbb Q}\arrow[u,dashed]\arrow[ur]\&
\end{tikzcd}}
\]
,
(which \deleted{\unsure{we}} surely does not exist) --- so
the obstruction is identified as the extension of $\Gamma_{\mathbb Q}$ by $\mu$,
inverse image of the extension $\mathcal{E}''$.
\margin{This is obvious in \unsure{effect|fact}, \unsure{that we|since} :
$(\Gamma_{\mathbb Q})_{\mathrm{ab}}\simeq\widehat{\mathbb Z}^{*}
\supset\mathbb Z_3^{*}$, now
$(\mathbb Z_3^{*})_4\simeq\mathbb Z/2\mathbb Z$
(\unsure{so already on} $\mathbb Z_3^{*}$ \unsure{it does not lift}\ldots}
(One should look at why, in

\margin{630}

\scanpage{632}

classical terms, this obstruction is \unsure{split}, when
one passes to the subgroup $\Gamma_{\mathbb Q}^{\prime\prime\prime}$ corresponding to the
field $\mathbb Q(\sqrt{3}))$

Thus, the $\mathbb Q$-rational point under consideration of
$M_{1,1}'$ does not lift to a $\mathbb Q$-rational point of
$M_{1,1}$ --- On the other hand, its image $\inserted{0}$ in
$\mathbb E^1\simeq U_{0,3}/\mathcal G_3$ must indeed
lift --- i.e. there must indeed exist a
$\widetilde{\mathcal G}_3$-torsor $\widetilde P$, and a
$\widetilde{\mathcal G}_3$-morphism
\deleted{$\unsure{P}$}

\[
\diagram{\begin{tikzcd}[ampersand replacement=\&,column sep=1.6em,row sep=0.1em]
\widetilde P\arrow[r]\&U_{0,3}\&\text{i.e.}\&P\arrow[r]\arrow[d,equal]\&U_{0,3}\\
\&\&\&\widetilde P/\mu\&
\end{tikzcd}}
\]
.

Let us consider indeed the \inserted{sub-}scheme
\deleted{of $U_{0,3}$, inverse image}
\[
\deleted{\unsure{Q=P}}\qquad
Q=U_{0,3}\mathbin{\times}_{\mathbb E^1}\{0\}
\]
the inverse image of $0$, i.e.
\[
P_0=Q_{\mathrm{red}}\ \deleted{\unsure{\simeq}}
\]
it is a subscheme which, over a field
$K\supset\mathbb Q(j)$, \deleted{splits into} is a sum of
two points, and which over $\mathbb Q$ is connected
--- it is a scheme \inserted{canonical\textsuperscript{ly}} isomorphic to
$\operatorname{Spec}\mathbb Q(-\bar j)$. It is stable under $\mathcal G_3$,
which operates on it via

\[
\diagram{\begin{tikzcd}[ampersand replacement=\&]
\mathcal G_3\arrow[r,"\mathrm{sg}"]\&\mu
\end{tikzcd}}
\]

and the Galois operation of $\mu$ on the $\mu_{\mathbb Q}$-torsor $P_0$.

\margin{631}

\scanpage{633}
\mspage{620}

Thus we have the diagram

\[
\diagram{\begin{tikzcd}[ampersand replacement=\&,column sep=1.8em,row sep=2em]
\&\widetilde{\mathcal G}_3\arrow[r,"\mathrm{can}"]\&\mathcal G_3\arrow[r,"\mathrm{sg}"]\&\mu\\
(31)\&\&\Gamma_{\mathbb Q}\arrow[ur]\&
\end{tikzcd}}
\]

\margin{\uline{N.B.} the kernel of
$\widetilde{\mathcal G}_3\longrightarrow\mu$ is a
\inserted{central} extension of $\mathbb Z/3\mathbb Z$
by $\mu=\mathbb Z/2\mathbb Z$, it
is therefore \unsure{trivial} and
isomorphic to $\mathbb Z/6\mathbb Z$
(kernel generated by $\rho$).}

We thus see that the category of liftings
of the $\mathbb Q$-rational point $0$ of $\mathbb E^1$ to a
$\mathbb Q$-rational point of $M_{1,1}'$ resp. $M_{1,1}$, is
equivalent to that of $\mathcal G_3$-torsors \inserted{$P$} resp\unsure{.}
of $\widetilde{\mathcal G}_3$-torsors \inserted{$\widetilde P$} over $\mathbb Q$, which lift the
$\mu$-torsor $P_0$. We had chosen a
$P=T$ in a \unsure{particularly} \unsure{trivial|trivial} way,
using the splitting

\[
\diagram{\begin{tikzcd}[ampersand replacement=\&,column sep=1.5em]
\mu\arrow[r,"\sim"]\&\{1,\dot\sigma_{\infty}\}\arrow[r,hook]\&\mathcal G_3
\end{tikzcd}}
\]

of $\mathcal G_3\longrightarrow\mu$ --- but the torsor $T$ does not
lift to a $\widetilde T$. The isomorphism classes of points of $M_{1,1}'$ resp. of $M_{1,1}$ over
the point $0$ of $\mathbb E^1$, \unsure{correspond} likewise to the
conjugacy classes \inserted{(modulo kernels \inserted{$\mathcal G_3^+,\widetilde{\mathcal G}_3^+$} of $\mathcal G_3\longrightarrow\mu$ and $\widetilde{\mathcal G}_3\longrightarrow\mu$)} of liftings of
$\Gamma_{\mathbb Q}\longrightarrow\mu$ to $\Gamma_{\mathbb Q}\longrightarrow\mathcal G_3$ resp. to $\Gamma_{\mathbb Q}\longrightarrow\widetilde{\mathcal G}_3$.

\deleted{\unsure{From there}} Let us consider a lifting to

\[
\diagram{\begin{tikzcd}[ampersand replacement=\&]
\Gamma_{\mathbb Q}\arrow[r]\&\mathcal G_3
\end{tikzcd}}
\]
,
its image is a subgroup of $\mathcal G_3$, which is either

\margin{632}

\scanpage{634}

$\{1,\dot\sigma_{\infty}\}$ or one of its conjugates,
$\{1,\dot\sigma_0\}$, $\{1,\unsure{\bar\sigma_1|\dot\sigma_1}\}$,
\uline{or} it is the whole of $\mathcal G_3$. In the first
case, we have seen that it does not lift
to $\widetilde{\mathcal G}_3$. So to be able to lift,
one must take an \uline{epimorphism} $\Gamma_{\mathbb Q}\longrightarrow\mathcal G_3$ ---
so we leave the abelian setting over
$\mathbb Q$, to enter the domain of
abelian extensions of $\mathbb Q(j)$.

To examine in terms of Galois
groups whether there exists such a lifting \deleted{which}
$\Gamma_{\mathbb Q}\longrightarrow\widetilde{\mathcal G}_3$, and \deleted{\unsure{which}}
how to obtain them,
appears as a pleasant
and \unsure{delectable} exercise in class field
theory, which I am not going to pursue
for the moment, so as not to make
this digression last forever. Returning later
to this point, I would need at the same
time to make explicit the relation between the
\unsure{choices|points} of a lifting of the homomorphism from $\Gamma_{\mathbb Q}$,
and splittings of the extension $\mathcal M_{0,3}$ of $\Gamma_{\mathbb Q}$
by $\unsure{\mathcal G_{0,3}^{+\wedge}|\widetilde{\mathcal G}_{0,3}^{+\wedge}}$,
or $\mathcal M$ of $\Gamma_{\mathbb Q}$ by
$\mathcal G^{+\wedge}\simeq\mathrm{SL}(2,\mathbb Z)^\wedge$.

\margin{633}

\scanpage{635}
\mspage{621}

Let us return to the subgroups \deleted{$\unsure{\mathcal N_{\rho}^{?}}$ of}

\[
\diagram{
\begin{tikzcd}[ampersand replacement=\&,column sep=3em,row sep=2em]
\mathcal N_{\rho}^{?}
  \arrow[r,phantom,"{\overset{3}{\unsure{\subset}}}"{description}]
\&\mathcal N_{\rho}^{*}\arrow[r,phantom,"{\subset}"{description}]
\&\mathcal M\\
Z_{\rho,0}^{!}=\mathcal M'(-\bar j)
  \arrow[u,phantom,"{4\,\cup}"{description}]
\&\&
\end{tikzcd}
}
\]

containing \unsure{the|the} partial \deleted{section}\inserted{\unsure{subgroup|subgroups}}
$\mathcal M(-\bar j)\simeq\widetilde\Gamma_{\mathbb Q}'$.

Let
\[
\tag{32}
\begin{aligned}
\mathcal M^0(-\bar j)&\overset{2}{\subset}\mathcal M'(-\bar j)\\
\underset{\unsure{\mathrm{dfn}|\mathrm{def}}}{=}&\operatorname{Ker}
\bigl(\mathcal M'(-\bar j)\xrightarrow{\varepsilon_2}\mu\bigr)
=\deleted{\unsure{\mathcal N_{\rho}^{?}|\mathcal M}}\,Z_{\rho}
\cap\deleted{\unsure{\operatorname{Ker}\varepsilon_2}}\,
\mathcal M_0^{\prime\prime!}\\
&=\mathcal N_{\rho}^{*}\cap\mathcal M_0^{0!}
\end{aligned}
\]
which is therefore a partial section \unsure{up|up} above
\[
\tag{33}
\begin{aligned}
\widetilde\Gamma_{\mathbb Q}^0
&=\operatorname{Ker}\bigl(\widetilde\Gamma_{\mathbb Q}
\xrightarrow{(\varepsilon_2,\varepsilon_3)}\mu\times\mu\bigr)
\deleted{\unsure{\ne|=}\widetilde\Gamma_{\mathbb Q}}\\
&=\operatorname{Ker}\bigl(\widetilde\Gamma_{\mathbb Q}
\longrightarrow\widehat{\mathbb Z}^{*}\longrightarrow
(\mathbb Z/12\mathbb Z)^{*}\bigr)
\end{aligned}
\]

\deleted{Let us consider} Modulo a verification
which I am going to make below, $\tau'$ normalizes
\deleted{(34)}
$\mathcal M^0(-\bar j)$, let us consider \unsure{the|the}
\[
\tag{34}
\mathcal M'''(-\bar j)\deleted{\unsure{\mathcal N_{\rho,0}^{?}}}
\overset{\unsure{\mathrm{dfn}|\mathrm{def}}}{=}
\underset{\text{(semidirect)}}{\mu_{\tau'}\mathbin{.}\mathcal M^0(-\bar j)}
\qquad\text{where }\mu_{\tau'}=\{1,\tau'\}
\]
it is a subgroup of index $4$ of $\mathcal N_{\rho}^{?}$,
and we now have
\[
\tag{35}
\deleted{\unsure{\mathcal N_{\rho,0}^{?}}}\quad
\mathcal M'''(-\bar j)\xrightarrow{\sim}\widetilde\Gamma_{\mathbb Q}''',
\]
it is a partial section \unsure{up|up} above
$\widetilde\Gamma_{\mathbb Q}'''$, whence isomorphisms

\margin{634}

\scanpage{636}

\deleted{
\[
\tag{36}
\left\{
\begin{aligned}
\mathcal M'''&\simeq\mathcal M'''(-\bar j)\mathbin{.}\mathcal G^{+\wedge}\\
\deleted{\unsure{\mathcal M'''}}\,\mathcal N_{\rho}^{*}
&\simeq\mathcal M'''(-\bar j)\mathbin{.}L_{\rho}.
\end{aligned}
\right.
\]

We have the summary diagram
of inclusions of subgroups of $\mathcal N_{\rho}^{*}$

\[
\diagram{
\begin{tikzcd}[ampersand replacement=\&,column sep=2.5em,row sep=1.5em]
\&\&\&Z_{\rho}\arrow[dr,hook,"2"]\&\\
\&\&\begin{gathered}\mathcal M'(-\bar j)\\=Z_{\rho,0}^{!}\end{gathered}
  \arrow[ur,hook,"6"]\arrow[dr,hook,"4"]
\&\&\mathcal N_{\rho}^{*}\\
(37)\&\mathcal M^0(-\bar j)\arrow[ur,hook,"2"]\arrow[dr,hook,"2"']
\&\&\mathcal N_{\rho}^{?}\arrow[ur,hook,"3"']\&\\
\&\&\mathcal M'''(-\bar j)\arrow[ur,hook,"4"]\&\&
\end{tikzcd}
}
\]


Let us note that the inclusion
\[
\tag{38}
\underset{\substack{\simeq\\\unsure{\mathbb D_6|\mathbb D_{\sigma}}}}{\mathbb D_{\rho}}
\overset{\mathrm{dfn}}{=}
\underset{\simeq\mathbb Z/2\mathbb Z}{\uline{\mu_{\tau'}}}
\mathbin{.}
\underset{=\mathbb Z/6\mathbb Z}{\uline{L_{\rho}}}
\subset\mathcal N_{\rho}^{*}
\]
induces an isomorphism of
\[
\tag{39}
\mathbb D_{\rho}\xrightarrow{\sim}\mathcal N_{\rho}^{*}/
\]
}

\margin{635}

\scanpage{637}
\mspage{622}

\[
\tag{36}
\left\{
\begin{aligned}
\mathcal{M}'''&\simeq \mathcal{M}'''(-\bar j)\mathbin{.}\mathcal{G}^{+\wedge}\\
\mathcal{N}_{\rho}^{*}&\simeq \mathcal{M}'''(-\bar j)\mathbin{.}L_{\rho}.
\end{aligned}
\right.
\]

Summary diagram of the subgroups of $\mathcal{N}_{\rho}^{*}$
obtained so far:


\[
\diagram{
\begin{tikzcd}[ampersand replacement=\&,column sep=1.05em,row sep=2.8em]
\&\& \mathcal{M}^{0}(-\bar j)
  \arrow[rr,hook,"2"]
  \arrow[dl,hook,"2"']
\&\& \begin{gathered}\mathcal{M}'(-\bar j)\\=Z_{\rho,0}^{!}\end{gathered}
  \arrow[rr,hook,"2"]
  \arrow[dl,hook,"2"']
  \arrow[dr,no head,draw=gray,
    "{\unsure{\mathcal{M}'(-\bar j)\times\mu}}"{description,pos=.5},
    "{\unsure{\searrow}}"{pos=.93}]
\&\& \mathcal{M}'(-\bar j)\times\mu
  \arrow[rr,hook,"3"]
  \arrow[dl,hook,"2"']
\&\& Z_{\rho}\arrow[dl,hook,"2"']\\
(37)\&\mathcal{M}'''(-\bar j)\arrow[rr,hook,"2"]
\&\& \mathcal{M}'''(-\bar j)\times\mu\arrow[rr,hook,"2"]
\&\& \mathcal{N}_{\rho}^{?}\arrow[rr,hook,"3"]
\&\& \mathcal{N}_{\rho}^{*}
\&\&\\
\&\&\&\&1
  \arrow[uu,hook]
  \arrow[uull,hook,start anchor=north west,end anchor=south west]
  \arrow[rr,hook,"2"]
  \arrow[dl,hook,"2"']
\&\&\mu\arrow[uu,hook]\arrow[rr,hook,"3"]\arrow[dl,hook,"2"']
\&\&L_{\rho}\arrow[uu,hook]\arrow[dl,hook,"2"']\\
\&\&\&\mu_{\tau'}\arrow[uu,hook]\arrow[uull,hook]\arrow[rr,hook,"2"]
\&\&\mathbb{D}_{\tau'}\arrow[uu,hook]\arrow[rr,hook,"3"]
\&\&\mathbb{D}_{\rho}\arrow[uu,hook]
\end{tikzcd}
}
\]
,

where we are here interested \unsure{above all} in the
``upper face'', \unsure{containing|concerning} subgroups of finite index in
$\mathcal{N}_{\rho}^{*}$, intermediate between it and
$\unsure{\mathcal{M}_{0}(-\bar j)|\mathcal{M}^{0}(-\bar j)}$
(of index $24$ in $\mathcal{N}_{\rho}^{\unsure{*}}$).
I am going to examine the question of invariance of these subgroups
in one Another, and the structure of the quotient groups.

I am going to redo a diagram which \unsure{really|better} shows the
\deleted{\unsure{groups|relationships}} relations between the groups in question.

\margin{636}